\phantomsection
\section*{\hfill РЕФЕРАТ \hfill}
\markboth{РЕФЕРАТ}{}

Расчетно-пояснительная записка 64 с., 8 рис., 21 табл., 32 источника, 1~приложение.

НЕЙРОННЫЙ МАШИННЫЙ ПЕРЕВОД, ТРАНСФОРМЕР, MAMBA, ENCODER-DECODER, РУССКО-АНГЛИЙСКИЙ ПЕРЕВОД, СРАВНЕНИЕ АРХИТЕКТУР

Объектом исследования являются архитектуры нейронных сетей для машинного перевода: LSTM, Transformer, Mamba, а также encoder-decoder и decoder-only схемы организации моделей.

Цель работы --- сравнительный анализ современных архитектурных и обучающих решений для NMT на примере русско-английского направления.

В работе рассмотрена постановка задачи NMT, методы токенизации, декодирования и оценки качества. Описаны архитектуры LSTM, Transformer, Mamba, а также encoder-decoder и decoder-only постановки. Подготовлены параллельные корпуса трёх размеров, выполнена их фильтрация и очистка. Проведено сравнение архитектур, позиционных кодирований, токенизаторов и гиперпараметров обучения. Выполнена проверка GRPO-дообучения.

Сравнение показало преимущество Transformer encoder-decoder с RoPE; гибридная Mamba encoder-decoder близка по качеству, decoder-only уступают. Качество определяется сочетанием архитектуры и режима обучения. Итоговая двуязычная модель с примерно 230 млн параметров достигла 32{,}25 BLEU, 60{,}34 chrF и 0{,}8519 COMET на тестовом наборе.

\newpage
